\documentclass[11pt]{article}
\usepackage{times}
\usepackage{fullpage}
\usepackage{url}
\usepackage{hyperref}
\usepackage{fancyhdr}
\usepackage{graphicx}
\usepackage{enumitem}
\usepackage{subcaption}
\usepackage{amsmath, amsfonts, amsthm, fouriernc}
\usepackage{IEEEtrantools}
\usepackage{parskip}
\usepackage{booktabs}

\pagestyle{fancy}
\addtolength{\headheight}{2em}
\addtolength{\headsep}{1.5em}
\lhead{\large{ME 2801}}

\usepackage{sectsty}
\sectionfont{\fontsize{12}{8}\selectfont}

% Auto-numbered exercise headings: \exercise{Title} -> "Exercise N --- Title"
\newcounter{exercise}
\newcommand{\exercise}[1]{%
  \refstepcounter{exercise}%
  \section*{Exercise \theexercise\ --- #1}%
}

\begin{document}

\begin{center}
  \Large{\bf{Assignment: Steady-State Error}}
\end{center}

This assignment exercises the steady-state error analysis from the chapter.
All computations assume \emph{unity negative feedback} and a \emph{stable}
closed loop unless stated otherwise.  

% ---------------------------------------------------------------
\exercise{Three Test Inputs}

The unity-feedback system has open-loop transfer function
\[
  G(s) = \frac{600}{(s+10)(s^2 + 4s + 25)}.
\]
For example, $G(s)$ might represent a spacecraft attitude-control loop:
a reaction-wheel actuator (the real pole at $s = -10$) driving a rigid
spacecraft bus coupled to a lightly-damped flexible solar panel (the
complex pole pair) --- recall from the chapter that $G(s)$ here is the
\emph{combined} open-loop $C(s)\,P(s)$, not the plant alone.

\begin{enumerate}[label=(\alph*)]
  \item Identify the system Type by inspection of $G(s)$.
  \item Find the steady-state error for the inputs $u(t)$, $t\,u(t)$, and
    $\tfrac{1}{2}\,t^2 u(t)$.  (The factor of $\tfrac12$ on the parabola is
    just a bookkeeping convenience: it gives $\mathcal{L}\{\tfrac12 t^2\} =
    1/s^3$, which keeps the algebra clean.)
  \item In MATLAB, verify your step-input answer in (b) two ways:
    \begin{itemize}
      \item Calculate: compute the closed-loop transfer function and
        evaluate its DC gain.
      \item Simulate: generate the step response and check the final value
        of the error signal.
    \end{itemize}
  \item Two of your answers in (b) are infinite.  In one sentence, explain
    what physical fact about $G(s)$ makes that so.
\end{enumerate}

\iffalse
% ---------------------------------------------------------------
\exercise{Speed (Transient) and Accuracy (Steady-State) Together}
% (Hidden via \iffalse...\fi --- kept in source for reference.)

The unity-feedback system has open-loop transfer function
\[
  G(s) = \frac{4000}{s(s+60)}.
\]

For each computational part below, work the answer first \emph{analytically
(by hand)}, then verify in MATLAB.

\begin{enumerate}[label=(\alph*)]
  \item Find the percent overshoot of the closed-loop unit-step response,
    by hand.  Verify with \texttt{stepinfo(feedback(G,1))}.
  \item Find the 2\,\% settling time, by hand.  Verify with the same
    \texttt{stepinfo} output.
  \item Find the steady-state error for the inputs $5u(t)$, $5t\,u(t)$, and
    $5t^2 u(t)$, by hand.  Verify each response with MATLAB.
  \item Suppose you want to reduce the ramp error in (c) by raising the
    loop gain.  In one or two sentences, explain why this might be a bad
    idea --- referencing your answer to (a).
\end{enumerate}
\fi

% ---------------------------------------------------------------
\exercise{Design a P-Only Velocity Tracking Compensator}

You will design a proportional-only controller for a position-control
plant.  The plant is
\[
  G_p(s) = \frac{s+5}{s\,(s+8)(s+12)},
\]
and the controller is $C(s) = K_p$.  The system is closed with unity
feedback.

For example, $G_p(s)$ might model a machine-tool feed drive, where
accurate constant feed velocity is required: the integrator $1/s$ converts the motor's angular velocity
to linear position, the real poles are the motor's electrical
and mechanical time constants, and the zero comes from a tachometer-feedback
element built into the plant.

Keep track of which transfer function you are working
with at each step: the \emph{open-loop, forward path} or the \emph{closed loop}.

\begin{enumerate}[label=(\alph*)]
  \item What is the system Type?
  \item Find the value of $K_p$ such that the steady-state error to a unit
    ramp input is 4\,\%.
  \item Use MATLAB to compute the closed-loop poles for your value of
    $K_p$.  Are they all in the left half-plane?
\end{enumerate}

\iffalse
% ---------------------------------------------------------------
\exercise{Fill the Table}
% (Hidden via \iffalse...\fi --- kept in source for reference.)

For each of the two open-loop transfer functions below in unity feedback,
fill in the steady-state error: a number, $0$, or $\infty$.
\[
  G_a(s) = \frac{20}{(s+2)(s+10)},
  \qquad
  G_b(s) = \frac{40}{s\,(s+8)}.
\]

\begin{enumerate}[label=(\alph*)]
  \item Identify the Type of each plant before computing.
  \item Fill in the six cells.
  \begin{center}
  \renewcommand{\arraystretch}{1.5}
  \begin{tabular}{l|ccc}
    \toprule
    & $r(t) = u(t)$ & $r(t) = t\,u(t)$ & $r(t) = \tfrac12 t^2 u(t)$ \\
    \midrule
    $G_a(s)$ & \rule{2.5cm}{0pt} & \rule{2.5cm}{0pt} & \rule{2.5cm}{0pt} \\
    $G_b(s)$ & & & \\
    \bottomrule
  \end{tabular}
  \end{center}
  \item In one sentence, what is the structural difference between $G_a$ and
    $G_b$ that explains the entire pattern in your table?
\end{enumerate}
\fi

% ---------------------------------------------------------------
\exercise{USV Surge: Predict Before You Simulate}

The USV surge (forward-speed) plant, identified from the lab data, is
\[
  G_\mathrm{surge}(s) = \frac{5.1}{1.3\,s + 1}.
\]
For each controller architecture below, find the steady-state error to a
unit-step (1\,m/s) speed command.  Use analytical methods, MATLAB, or both.

\begin{enumerate}[label=(\alph*)]
  \item \textbf{P only}: $K_P = 0.5$.
  \item \textbf{PI}: $K_P = 0.5$, $K_I = 0.4$.
  \item \textbf{P + feedforward}: $K_P = 0.5$, $K_{\!f\!f} = 1/K_\mathrm{plant}$.
  \item \textbf{PI + feedforward}: same $K_P, K_I$ as (b), and $K_{\!f\!f}$ as (c).
  \item Use the \emph{adaptation argument} from the chapter (not the
    steady-state error formula) to explain in one or two sentences why the
    PI controller in (b) drives the error to zero.
\end{enumerate}

% ---------------------------------------------------------------
\exercise{USV Yaw: Ramp Tracking Design}

The USV's yaw-rate plant, identified from the lab data, is
\[
  G_\mathrm{yaw\text{-}rate}(s) = \frac{0.82}{0.61\,s + 1}.
\]
We want to control \emph{heading angle}, not yaw rate, to follow a course.  Heading is the
integral of yaw rate.  A proportional controller $C(s) = K_r$ is used in a
unity-feedback heading loop.  The heading reference (the command) ramps at $5^\circ$/s
(a constant turn-rate command).

\begin{enumerate}[label=(\alph*)]
  \item Write the plant transfer function from rudder input to
    \emph{heading angle}, $G_r(s)$.  (Hint: heading is the integral of yaw
    rate, so $G_r(s)$ picks up an extra factor relative to the yaw-rate
    plant above.)
  \item What is the system Type of the closed loop?  Justify in one phrase.
  \item Write
    the steady-state heading error in degrees as a function of $K_r$.
  \item The design spec is an absolute heading error $\le 1^\circ$.  Find the minimum
    $K_r$ that meets the spec.
  \item Verify your $K_r$ in MATLAB by simulating the closed-loop ramp
    response and checking the steady-state error.
\end{enumerate}

\end{document}
